\chapter{Product Measures and Fubini Theory}
\label{ch:iii05}

Product measures combine measurable coordinate spaces, while Tonelli and Fubini convert multidimensional integration into iterated integration. Absolute integrability is the key condition that permits changing order safely.

\section*{Learning goals}
After this chapter, the reader should be able to:
\begin{itemize}
\item construct product sigma-algebras and product measures in standard sigma-finite settings;
\item apply Tonelli's theorem to nonnegative functions and Fubini's theorem to integrable functions with the correct hypotheses;
\item justify changing the order of integration by checking absolute integrability or nonnegativity;
\item compute iterated integrals on rectangles and simple product domains;
\item identify measurable sections and use them to interpret product-measure statements;
\item construct examples where changing integration order fails because the Fubini hypotheses are not satisfied;
\end{itemize}
\section*{Conceptual roadmap}
\[
\boxed{\text{structure}\;\longrightarrow\;\text{estimate}\;\longrightarrow\;\text{limit or representation}\;\longrightarrow\;\text{application}.}
\]

\paragraph{Mathematical setting.}
Throughout this chapter $(X,\mathcal A,\mu)$ and $(Y,\mathcal B,\nu)$ are sigma-finite measure spaces, and product measurability means $\mathcal A\otimes\mathcal B$, before completion. The rectangle convention is $0\cdot\infty=0$. For the completed product, measurable representatives in the uncompleted product give the same integrals; section assertions then hold a.e. in the completed factor spaces.

\section{Product sigma-algebra}
The product sigma-algebra is generated by measurable rectangles $A\times B$.

\section{Product measures}
On sigma-finite spaces, the product measure is characterized by $(\mu\times\nu)(A\times B)=\mu(A)\nu(B)$ on rectangles.

\section{Sections of sets}
For $E\in\mathcal A\otimes\mathcal B$, the section $E_x=\{y:(x,y)\in E\}$ belongs to $\mathcal B$ for every $x$: sets with measurable sections form a sigma-algebra containing the rectangles.

\section{Sections of functions}
Fixing one variable in a measurable function produces measurable sections, allowing iterated integration.

\section{Tonelli theorem}
Nonnegative measurable functions may be integrated in either order, with the common value possibly infinite.


% BEGIN VOL03-EXPANSION III05-example-01
\begin{example}[Tonelli for a nonnegative product]\label{ex:iii05-expand-01}
For \(f(x,y)=e^{-x-y}\) on \([0,\infty)^2\), Tonelli applies immediately
because \(f\ge0\). Either iterated integral equals
\[
\left(\int_0^\infty e^{-x}dx\right)
\left(\int_0^\infty e^{-y}dy\right)=1.
\]
Nonnegativity legitimizes the interchange even before finiteness is known.
\end{example}
% END VOL03-EXPANSION III05-example-01
\section{Fubini theorem}
Absolutely integrable functions have integrable sections almost everywhere and finite iterated integrals equal to the product integral.


% BEGIN VOL03-EXPANSION III05-example-02
\begin{example}[Fubini after checking absolute integrability]\label{ex:iii05-expand-02}
On \([0,1]^2\), \(f(x,y)=x-y\) satisfies \(|f|\le1\), so \(f\in L^1\).
Fubini therefore applies and either integration order gives zero. The
absolute-integrability check is the structural step that distinguishes a legal
change of order from a formal rearrangement.
\end{example}
% END VOL03-EXPANSION III05-example-02
\section{Changing order}
Absolute integrability prevents order-dependent cancellation.


% BEGIN VOL03-EXPANSION III05-example-03
\begin{example}[Changing order needs a theorem]\label{ex:iii05-expand-03}
For a signed function with divergent absolute integral, positive and negative
mass may cancel differently in different iterated orders. Thus the diagnostic
is to inspect the integral of \(|f|\). If it is infinite, Fubini cannot be
invoked; each iterated integral must be analyzed independently.
\end{example}
% END VOL03-EXPANSION III05-example-03
\section{Applications}
Areas, convolution kernels, expectations, and integral operators all rely on product integration.

\section{Core structural results}

\begin{theorem}[Tonelli theorem]\label{thm:iii05-01}
For nonnegative measurable $f$ on a sigma-finite product space, the product integral equals both iterated integrals, possibly with value $\infty$.
\begin{proof}
Proof sketch (standard product-measure theorem). On finite-measure rectangles, a pi-lambda argument extends the section-integral identity from rectangles to all product-measurable sets; sigma-finite exhaustions extend it to the full product. Only then pass from indicators to nonnegative simple functions and use MCT. The measure-extension and pi-lambda theorems used in this step are standard theorems whose full proofs are deferred.
\end{proof}
\end{theorem}

\begin{theorem}[Fubini theorem]\label{thm:iii05-02}
If $f\in L^1(\mu\times\nu)$, then almost all sections are integrable and the iterated integrals equal $\int f$.
\begin{proof}
Apply Tonelli to $|f|$ to obtain finite section integrals almost everywhere, then decompose $f$ into positive and negative parts.
\end{proof}
\end{theorem}

\begin{theorem}[Rectangle rule]\label{thm:iii05-03}
For measurable rectangles, $(\mu\times\nu)(A\times B)=\mu(A)\nu(B)$.
\begin{proof}
Standard theorem / proof sketch: the rectangle values extend additively to the algebra of finite disjoint unions of rectangles. The product-measure construction proves this is a premeasure; the measure-extension theorem and sigma-finiteness give existence and uniqueness. These construction theorems are used here with their full proofs deferred.
\end{proof}
\end{theorem}

\section{Worked examples}

\begin{example}[Product sigma-algebra diagnostic]\label{ex:iii05-worked-01}
Give a rigorous worked analysis of Product sigma-algebra in the setting of this chapter. State the decisive definition or hypothesis and derive the central conclusion. The product sigma-algebra is generated by measurable rectangles $A\times B$. The decisive point is to use the chapter definition at the correct countable, norm, transform, or distributional level rather than rely on pointwise intuition alone.
\end{example}

\begin{example}[Product measures diagnostic]\label{ex:iii05-worked-02}
Give a rigorous worked analysis of Product measures in the setting of this chapter. State the decisive definition or hypothesis and derive the central conclusion. On sigma-finite spaces, the product measure is characterized by $(\mu\times\nu)(A\times B)=\mu(A)\nu(B)$ on rectangles. The decisive point is to use the chapter definition at the correct countable, norm, transform, or distributional level rather than rely on pointwise intuition alone.
\end{example}

\begin{example}[Sections of sets diagnostic]\label{ex:iii05-worked-03}
Give a rigorous worked analysis of Sections of sets in the setting of this chapter. State the decisive definition or hypothesis and derive the central conclusion. For $E\in\mathcal A\otimes\mathcal B$, the section $E_x=\{y:(x,y)\in E\}$ belongs to $\mathcal B$ for every $x$: sets with measurable sections form a sigma-algebra containing the rectangles. The decisive point is to use the chapter definition at the correct countable, norm, transform, or distributional level rather than rely on pointwise intuition alone.
\end{example}

\begin{example}[Sections of functions diagnostic]\label{ex:iii05-worked-04}
Give a rigorous worked analysis of Sections of functions in the setting of this chapter. State the decisive definition or hypothesis and derive the central conclusion. Fixing one variable in a measurable function produces measurable sections, allowing iterated integration. The decisive point is to use the chapter definition at the correct countable, norm, transform, or distributional level rather than rely on pointwise intuition alone.
\end{example}

\section{Solved dossiers}
Each dossier is a canonical solved problem. The provenance ledger records which retained corpus rules guided its topic and which dossiers were added freshly to complete the chapter.

\begin{problem}[Product sigma-algebra diagnostic]\label{prob:iii05-dossier-01}
Give a rigorous worked analysis of Product sigma-algebra in the setting of this chapter. State the decisive definition or hypothesis and derive the central conclusion.
\end{problem}
\begin{solution}
The product sigma-algebra is generated by measurable rectangles $A\times B$. The decisive point is to use the chapter definition at the correct countable, norm, transform, or distributional level rather than rely on pointwise intuition alone.
\end{solution}

\begin{problem}[Product measures diagnostic]\label{prob:iii05-dossier-02}
Give a rigorous worked analysis of Product measures in the setting of this chapter. State the decisive definition or hypothesis and derive the central conclusion.
\end{problem}
\begin{solution}
On sigma-finite spaces, the product measure is characterized by $(\mu\times\nu)(A\times B)=\mu(A)\nu(B)$ on rectangles. The decisive point is to use the chapter definition at the correct countable, norm, transform, or distributional level rather than rely on pointwise intuition alone.
\end{solution}

\begin{problem}[Sections of sets diagnostic]\label{prob:iii05-dossier-03}
Give a rigorous worked analysis of Sections of sets in the setting of this chapter. State the decisive definition or hypothesis and derive the central conclusion.
\end{problem}
\begin{solution}
For $E\in\mathcal A\otimes\mathcal B$, the section $E_x=\{y:(x,y)\in E\}$ belongs to $\mathcal B$ for every $x$: sets with measurable sections form a sigma-algebra containing the rectangles. The decisive point is to use the chapter definition at the correct countable, norm, transform, or distributional level rather than rely on pointwise intuition alone.
\end{solution}

\begin{problem}[Sections of functions diagnostic]\label{prob:iii05-dossier-04}
Give a rigorous worked analysis of Sections of functions in the setting of this chapter. State the decisive definition or hypothesis and derive the central conclusion.
\end{problem}
\begin{solution}
Fixing one variable in a measurable function produces measurable sections, allowing iterated integration. The decisive point is to use the chapter definition at the correct countable, norm, transform, or distributional level rather than rely on pointwise intuition alone.
\end{solution}

\begin{problem}[Tonelli theorem diagnostic]\label{prob:iii05-dossier-05}
Give a rigorous worked analysis of Tonelli theorem in the setting of this chapter. State the decisive definition or hypothesis and derive the central conclusion.
\end{problem}
\begin{solution}
Nonnegative measurable functions may be integrated in either order, with the common value possibly infinite. The decisive point is to use the chapter definition at the correct countable, norm, transform, or distributional level rather than rely on pointwise intuition alone.
\end{solution}

\begin{problem}[Fubini theorem diagnostic]\label{prob:iii05-dossier-06}
Give a rigorous worked analysis of Fubini theorem in the setting of this chapter. State the decisive definition or hypothesis and derive the central conclusion.
\end{problem}
\begin{solution}
Absolutely integrable functions have integrable sections almost everywhere and finite iterated integrals equal to the product integral. The decisive point is to use the chapter definition at the correct countable, norm, transform, or distributional level rather than rely on pointwise intuition alone.
\end{solution}

\begin{problem}[Changing order diagnostic]\label{prob:iii05-dossier-07}
Give a rigorous worked analysis of Changing order in the setting of this chapter. State the decisive definition or hypothesis and derive the central conclusion.
\end{problem}
\begin{solution}
Absolute integrability prevents order-dependent cancellation. The decisive point is to use the chapter definition at the correct countable, norm, transform, or distributional level rather than rely on pointwise intuition alone.
\end{solution}

\begin{problem}[Applications diagnostic]\label{prob:iii05-dossier-08}
Give a rigorous worked analysis of Applications in the setting of this chapter. State the decisive definition or hypothesis and derive the central conclusion.
\end{problem}
\begin{solution}
Areas, convolution kernels, expectations, and integral operators all rely on product integration. The decisive point is to use the chapter definition at the correct countable, norm, transform, or distributional level rather than rely on pointwise intuition alone.
\end{solution}

\begin{problem}[Tonelli theorem proof dossier]\label{prob:iii05-dossier-09}
Explain the stated proof sketch using the explicitly deferred standard theorem inputs, and identify where the hypotheses enter: For nonnegative measurable $f$ on a sigma-finite product space, the product integral equals both iterated integrals, possibly with value $\infty$.
\end{problem}
\begin{solution}
Proof sketch (standard product-measure theorem). On finite-measure rectangles, a pi-lambda argument extends the section-integral identity from rectangles to all product-measurable sets; sigma-finite exhaustions extend it to the full product. Only then pass from indicators to nonnegative simple functions and use MCT. The measure-extension and pi-lambda theorems used in this step are standard theorems whose full proofs are deferred. This is a proof sketch with the scope stated above.
\end{solution}

\begin{problem}[Fubini theorem proof dossier]\label{prob:iii05-dossier-10}
Prove the following structural statement and identify the step where its main hypothesis is used: If $f\in L^1(\mu\times\nu)$, then almost all sections are integrable and the iterated integrals equal $\int f$.
\end{problem}
\begin{solution}
Apply Tonelli to $|f|$ to obtain finite section integrals almost everywhere, then decompose $f$ into positive and negative parts. This proof also explains why the stated hypothesis is structural rather than cosmetic.
\end{solution}

\begin{problem}[Rectangle rule proof dossier]\label{prob:iii05-dossier-11}
Explain the stated proof sketch using the explicitly deferred standard theorem inputs, and identify where the hypotheses enter: For measurable rectangles, $(\mu\times\nu)(A\times B)=\mu(A)\nu(B)$.
\end{problem}
\begin{solution}
Standard theorem / proof sketch: the rectangle values extend additively to the algebra of finite disjoint unions of rectangles. The product-measure construction proves this is a premeasure; the measure-extension theorem and sigma-finiteness give existence and uniqueness. These construction theorems are used here with their full proofs deferred. This is a proof sketch with the scope stated above.
\end{solution}

\begin{problem}[Chapter synthesis]\label{prob:iii05-dossier-12}
Connect the main definitions and structural theorems of this chapter into one proof strategy. Explain what object is controlled, what estimate or closure principle is used, and what conclusion becomes available.
\end{problem}
\begin{solution}
Product measures combine measurable coordinate spaces, while Tonelli and Fubini convert multidimensional integration into iterated integration. Absolute integrability is the key condition that permits changing order safely. The recurring strategy is to encode the object in the chapter's natural measurable, normed, Fourier, or distributional structure, establish the controlling estimate, and then pass to the desired limit or representation.
\end{solution}

\section{Exercises with complete solutions}
The shorter exercise layer remains separate from the solved dossiers.

\begin{exercise}\label{exr:iii05-01}
State and apply the central principle behind Product sigma-algebra. Give one immediate consequence in the setting of this chapter.
\end{exercise}
\begin{hint}
On rectangles, start from measurable rectangles and extend to the product sigma-algebra before using product measure.
\end{hint}
\begin{solution}
The product sigma-algebra is generated by measurable rectangles $A\times B$. As an immediate consequence, the same principle can be inserted into later convergence, approximation, transform, or weak-form arguments whenever its hypotheses are verified.
\end{solution}

\begin{exercise}\label{exr:iii05-02}
State and apply the central principle behind Product measures. Give one immediate consequence in the setting of this chapter.
\end{exercise}
\begin{hint}
Tonelli applies to nonnegative functions even when the integral is infinite; use it before worrying about absolute integrability.
\end{hint}
\begin{solution}
On sigma-finite spaces, the product measure is characterized by $(\mu\times\nu)(A\times B)=\mu(A)\nu(B)$ on rectangles. As an immediate consequence, the same principle can be inserted into later convergence, approximation, transform, or weak-form arguments whenever its hypotheses are verified.
\end{solution}

\begin{exercise}\label{exr:iii05-03}
State and apply the central principle behind Sections of sets. Give one immediate consequence in the setting of this chapter.
\end{exercise}
\begin{hint}
Fubini requires integrability; check \(\int |f|<\infty\) before changing the integration order.
\end{hint}
\begin{solution}
For $E\in\mathcal A\otimes\mathcal B$, the section $E_x=\{y:(x,y)\in E\}$ belongs to $\mathcal B$ for every $x$: sets with measurable sections form a sigma-algebra containing the rectangles. As an immediate consequence, the same principle can be inserted into later convergence, approximation, transform, or weak-form arguments whenever its hypotheses are verified.
\end{solution}

\begin{exercise}\label{exr:iii05-04}
State and apply the central principle behind Sections of functions. Give one immediate consequence in the setting of this chapter.
\end{exercise}
\begin{hint}
Compute the inner integral as a function of the outer variable and verify its measurability.
\end{hint}
\begin{solution}
Fixing one variable in a measurable function produces measurable sections, allowing iterated integration. As an immediate consequence, the same principle can be inserted into later convergence, approximation, transform, or weak-form arguments whenever its hypotheses are verified.
\end{solution}

\begin{exercise}\label{exr:iii05-05}
State and apply the central principle behind Tonelli theorem. Give one immediate consequence in the setting of this chapter.
\end{exercise}
\begin{hint}
For indicator functions, reduce Fubini/Tonelli to measuring sections of the underlying measurable set.
\end{hint}
\begin{solution}
Nonnegative measurable functions may be integrated in either order, with the common value possibly infinite. As an immediate consequence, the same principle can be inserted into later convergence, approximation, transform, or weak-form arguments whenever its hypotheses are verified.
\end{solution}

\begin{exercise}\label{exr:iii05-06}
State and apply the central principle behind Fubini theorem. Give one immediate consequence in the setting of this chapter.
\end{exercise}
\begin{hint}
Sigma-finiteness is structural in standard product-measure construction; identify the finite-measure pieces explicitly.
\end{hint}
\begin{solution}
Absolutely integrable functions have integrable sections almost everywhere and finite iterated integrals equal to the product integral. As an immediate consequence, the same principle can be inserted into later convergence, approximation, transform, or weak-form arguments whenever its hypotheses are verified.
\end{solution}

\begin{exercise}\label{exr:iii05-07}
State and apply the central principle behind Changing order. Give one immediate consequence in the setting of this chapter.
\end{exercise}
\begin{hint}
To expose failure of Fubini, look for conditionally convergent positive and negative parts with infinite total variation.
\end{hint}
\begin{solution}
Absolute integrability prevents order-dependent cancellation. As an immediate consequence, the same principle can be inserted into later convergence, approximation, transform, or weak-form arguments whenever its hypotheses are verified.
\end{solution}

\begin{exercise}\label{exr:iii05-08}
State and apply the central principle behind Applications. Give one immediate consequence in the setting of this chapter.
\end{exercise}
\begin{hint}
After swapping integrals, compare the two iterated expressions to the original product-space integral rather than only to each other.
\end{hint}
\begin{solution}
Areas, convolution kernels, expectations, and integral operators all rely on product integration. As an immediate consequence, the same principle can be inserted into later convergence, approximation, transform, or weak-form arguments whenever its hypotheses are verified.
\end{solution}


% BEGIN VOL03-EXPANSION III05-exercises-01
\section{Graded supplementary exercises}
These exercises extend the preserved foundational set and are graded by mathematical role.

\subsection*{Standard computations}

\begin{exercise}[Rectangle]\label{exr:iii05-expand-01}
Find the product measure of a rectangle with side measures 2 and 3.
\end{exercise}
\begin{hint}
Multiply.
\end{hint}
\begin{solution}
The measure is 6.
\end{solution}

\begin{exercise}[Separable integral]\label{exr:iii05-expand-02}
Integrate xy over the unit square.
\end{exercise}
\begin{hint}
Factor the iterated integral.
\end{hint}
\begin{solution}
The value is 1/4.
\end{solution}

\begin{exercise}[Section]\label{exr:iii05-expand-03}
For 0<y<x<1, describe the y-section at fixed x.
\end{exercise}
\begin{hint}
Fix x.
\end{hint}
\begin{solution}
It is (0,x) when 0<x<1, and empty otherwise.
\end{solution}

\begin{exercise}[Tonelli infinity]\label{exr:iii05-expand-04}
If a nonnegative iterated integral is infinite, what does Tonelli say about the product integral?
\end{exercise}
\begin{hint}
Use equality in the extended sense.
\end{hint}
\begin{solution}
It is infinite as well.
\end{solution}

\begin{exercise}[Fubini criterion]\label{exr:iii05-expand-05}
Name a standard sufficient condition for changing order for a signed function.
\end{exercise}
\begin{hint}
Check absolute integrability.
\end{hint}
\begin{solution}
Membership in L1 of the product space.
\end{solution}

\subsection*{Proofs}

\begin{exercise}[Rectangle rule]\label{exr:iii05-expand-06}
Explain the product rectangle rule.
\end{exercise}
\begin{hint}
It defines the product premeasure.
\end{hint}
\begin{solution}
Extension preserves mu(A)nu(B) on rectangles.
\end{solution}

\begin{exercise}[Tonelli proof]\label{exr:iii05-expand-07}
Outline Tonelli's extension steps.
\end{exercise}
\begin{hint}
Indicators, simple functions, then monotone convergence.
\end{hint}
\begin{solution}
First extend the rectangle identity to arbitrary measurable-set indicators by the pi-lambda and sigma-finite exhaustion steps in the theorem sketch. Then take finite positive sums and monotone limits.
\end{solution}

\begin{exercise}[Fubini from Tonelli]\label{exr:iii05-expand-08}
Derive Fubini from Tonelli.
\end{exercise}
\begin{hint}
Apply Tonelli to |f|, then split signs.
\end{hint}
\begin{solution}
Absolute integrability makes positive and negative parts finite enough for subtraction.
\end{solution}

\begin{exercise}[Section measurability]\label{exr:iii05-expand-09}
Why are measurable sections needed?
\end{exercise}
\begin{hint}
Iterated integrals require measurable one-variable functions.
\end{hint}
\begin{solution}
Product measurability supplies the section structure.
\end{solution}

\subsection*{Counterexamples and hypothesis testing}

\begin{exercise}[Conditional order]\label{exr:iii05-expand-10}
Why can changing order fail without absolute integrability?
\end{exercise}
\begin{hint}
Cancellation can depend on rearrangement.
\end{hint}
\begin{solution}
Positive and negative infinite masses can regroup differently.
\end{solution}

\begin{exercise}[Sigma-finite caveat]\label{exr:iii05-expand-11}
Why is sigma-finiteness prominent in product theory?
\end{exercise}
\begin{hint}
It supports extension uniqueness and standard Fubini/Tonelli formulations.
\end{hint}
\begin{solution}
Without it, product-measure behavior can be subtler.
\end{solution}

\begin{exercise}[Iterated existence]\label{exr:iii05-expand-12}
Why does existence of two iterated integrals not automatically prove Fubini?
\end{exercise}
\begin{hint}
Absolute integrability may still fail.
\end{hint}
\begin{solution}
For example, on $(0,1)^2$ set $f(x,y)=(x^2-y^2)/(x^2+y^2)^2$. Since $f=\partial_y[y/(x^2+y^2)]=-\partial_x[x/(x^2+y^2)]$, integration in $y$ then $x$ gives $\pi/4$, and in $x$ then $y$ gives $-\pi/4$. Polar coordinates show $\int|f|=\infty$ near the origin. Thus both finite iterated integrals can exist and disagree.
\end{solution}

\subsection*{Applications and investigations}

\begin{exercise}[Probability]\label{exr:iii05-expand-13}
How does Fubini appear in expectations of two variables?
\end{exercise}
\begin{hint}
Use a product probability space.
\end{hint}
\begin{solution}
Integrable joint functions can be averaged conditionally in either order.
\end{solution}

\begin{exercise}[Convolution theorem]\label{exr:iii05-expand-14}
Why is Fubini used in proving the convolution theorem?
\end{exercise}
\begin{hint}
The transform of a convolution is a double integral.
\end{hint}
\begin{solution}
Absolute integrability permits swapping and factoring the integrals.
\end{solution}

\subsection*{Challenge problems}

\begin{exercise}[Triangle domain]\label{exr:iii05-expand-15}
Reverse the order in the integral over 0<y<x<1.
\end{exercise}
\begin{hint}
Describe the same triangle from y first.
\end{hint}
\begin{solution}
It becomes y from 0 to 1 and x from y to 1.
\end{solution}

\begin{exercise}[A.e. section integrability]\label{exr:iii05-expand-16}
Why are sections of an L1 product function integrable almost everywhere?
\end{exercise}
\begin{hint}
Apply Tonelli to |f|.
\end{hint}
\begin{solution}
The function of the section integral is integrable, hence finite almost everywhere.
\end{solution}
% END VOL03-EXPANSION III05-exercises-01
\section*{Chapter summary}
The chapter now supplies a canonical theorem-and-dossier layer for subsequent measure, Fourier, distributional, Sobolev, and PDE arguments.
